% !TeX root = main.tex
\subsection{基于 BiMSeg 与 XGBoost 的枯叶率提取}
\label{sec:ryegrass-senescence}

图像枯叶率的估计取决于两个环节：能否准确确定黑麦草冠层的投影范围，以及能否在该范围内区分绿叶与枯叶。裸土等背景与枯叶的颜色可能相近，直接对整幅图像进行颜色分类容易将背景计入枯叶。为此，本研究首先利用 BiMSeg 提取黑麦草冠层，再利用 XGBoost 对掩膜内的像素进行分类，最后根据绿叶与枯叶的像素数计算图像枯叶率。

\subsubsection{绿叶与枯叶像素分类器构建}
\label{sec:ryegrass-xgboost-training}

从 100 张黑麦草 RGB 图像中人工点选可辨识的绿叶与枯叶像素，共获得 119\,731 个带类别标签的像素样本，其中绿叶标记为 0，枯叶标记为 1。标注像素、颜色特征与类别标签的对应关系及分类器训练流程见图~\ref{fig:ryegrass-xgboost-training}。
\begin{figure}[H]
  \centering
  \includegraphics[width=0.95\linewidth]{figures/xgboost_leaf_classification.pdf}
  \methodbicaption{黑麦草绿叶与枯叶像素分类器的训练流程}
    {Training workflow of the XGBoost classifier for green and senescent ryegrass leaf pixels}
    {fig:ryegrass-xgboost-training}
\end{figure}

每个标注点对应一个像素样本。以该像素的原始 RGB 值为基础，构建 14 维颜色特征向量，包括 \(R\)、\(G\)、\(B\)，HSV 空间中的 \(H\)、\(S\)，Lab 空间中的 \(a^*\)、\(b^*\)，以及灰度特征和 \(M\)、\(Y\)、\(C_b\)、\(C_r\)、\(I\)、\(Q\) 分量。训练集与验证集按 \(9:1\) 划分。训练时，以颜色特征向量为输入，以人工标注的类别为监督信号，拟合 XGBoost 二分类器。与基于单一颜色分量的固定阈值判别相比，该模型能够利用多个颜色分量的组合识别绿叶与枯叶，为冠层内的逐像素统计提供分类依据。

\subsubsection{绿叶与枯叶像素分类性能}
\label{sec:ryegrass-xgboost-performance}

采用总体准确率、绿叶 F1、枯叶 F1 和 ROC 曲线下面积（ROC AUC）评价分类器的判别性能。总体准确率反映两类像素的整体分类表现，类别 F1 综合反映相应类别的精确率与召回率，ROC AUC 则衡量模型在不同判别阈值下区分两类像素的能力。前期评估结果见图~\ref{fig:ryegrass-xgboost-performance}。
\begin{figure}[H]
  \centering
  \begin{tikzpicture}
    \begin{axis}[
      ybar,
      width=0.90\linewidth,
      height=5.6cm,
      xmin=-0.55, xmax=3.55,
      ymin=0, ymax=110,
      xtick={0,1,2,3},
      xticklabels={总体准确率,绿叶 F1,枯叶 F1,ROC AUC},
      ytick={0,20,40,60,80,100},
      ylabel={指标值（\%）},
      axis lines=left,
      tick style={black!65},
      tick label style={font=\small},
      label style={font=\small},
      ymajorgrids,
      grid style={black!12},
      bar width=23pt,
      point meta=y,
      nodes near coords={
        \pgfmathprintnumber[fixed,fixed zerofill,precision=2]{\pgfplotspointmeta}\%
      },
      every node near coord/.append style={font=\small,yshift=2pt},
    ]
      \addplot[bar shift=0pt,draw=blue!40!black,fill=blue!28!gray]
        coordinates {(0,94.05)};
      \addplot[bar shift=0pt,draw=green!30!black,fill=green!28!gray]
        coordinates {(1,96.13)};
      \addplot[bar shift=0pt,draw=orange!60!black,fill=orange!60!gray]
        coordinates {(2,87.12)};
      \addplot[bar shift=0pt,draw=blue!40!black,fill=blue!28!gray]
        coordinates {(3,97.92)};
    \end{axis}
  \end{tikzpicture}
  \methodbicaption{XGBoost 对黑麦草绿叶与枯叶像素的分类性能}
    {Performance of XGBoost in classifying green and senescent ryegrass leaf pixels}
    {fig:ryegrass-xgboost-performance}
  \par\smallskip
  \noindent\begin{minipage}{\linewidth}
    \normalfont\zihao{5}\raggedright
    \noindent 注：各指标均以百分数表示。F1 为精确率与召回率的调和平均，分别对绿叶和枯叶像素计算；ROC AUC 为受试者工作特征曲线下面积。\par
    \vspace{2pt}
    \noindent Note: All metrics are reported as percentages. F1 is the harmonic mean of precision and recall, computed separately for green and senescent leaf pixels. ROC AUC is the area under the receiver operating characteristic curve.\par
  \end{minipage}\par
\end{figure}

XGBoost 的总体准确率为 94.05\%，ROC AUC 为 0.9792，表明由颜色特征得到的分类得分能够较好地区分所评估的绿叶与枯叶像素。绿叶和枯叶的 F1 分数分别为 96.13\% 和 87.12\%，说明分类器对两类像素均具有识别能力，其中枯叶的判别表现低于绿叶。由于枯叶像素数直接进入枯叶率的分子，单独报告枯叶 F1 比仅报告总体准确率更能说明该分类器在表型提取中的作用。这些结果为采用 XGBoost 识别冠层内的枯叶组织提供了像素尺度的依据。

\subsubsection{联合推理与枯叶率计算}
\label{sec:ryegrass-senescence-inference}

推理时，首先由 BiMSeg 生成黑麦草冠层的前景掩膜。随后，从原始 RGB 图像中提取掩膜内像素的 14 维颜色特征，并由训练完成的 XGBoost 对这些像素逐一分类。预测标签映射回原图坐标后，得到绿叶与枯叶在冠层内的空间分布。不同背景下的联合推理结果见图~\ref{fig:ryegrass-senescence-results}。
\begin{figure}[H]
  \centering
  \includegraphics[width=0.95\linewidth]{figures/bimseg_xgboost_ryegrass_results.pdf}
  \methodbicaption{BiMSeg 与 XGBoost 对黑麦草绿叶与枯叶的联合识别结果}
    {Pixel-wise classification of green and senescent ryegrass leaves using BiMSeg and XGBoost}
    {fig:ryegrass-senescence-results}
  \par\smallskip
  {\raggedright\normalfont\zihao{5}
   注：a 为原始图像；b 为像素分类结果，绿色、黄色和黑色分别表示绿叶、枯叶和背景；c 为分类结果与原图的叠加图。每列对应同一样本。\par
   Note: a, original images; b, pixel classification maps, with green,
   yellow, and black indicating green leaves, senescent leaves, and background,
   respectively; c, classification overlays. Each column shows the same sample.\par}
\end{figure}

在所示样本中，底布与裸土在分类图中大部分被排除，两类预测区域主要沿可见叶片分布。叠加图进一步将分类结果与原图中的叶片位置对应起来，可用于检查颜色相近的背景是否被误计，以及细窄叶缘和绿枯交界处是否存在明显误分。视觉结果与像素分类评估分别从空间位置和类别判别两个方面支持该流程用于枯叶率估计。

设单幅图像中预测为绿叶和枯叶的像素数分别为 \(N_{\mathrm{green}}\) 和 \(N_{\mathrm{sen}}\)。当两者之和大于零时，图像枯叶率定义为
\begin{equation}
  R_{\mathrm{sen}} =
  \frac{N_{\mathrm{sen}}}
       {N_{\mathrm{green}}+N_{\mathrm{sen}}}
  \times 100\%.
  \label{eq:ryegrass-senescence-rate}
\end{equation}
式~\eqref{eq:ryegrass-senescence-rate} 表示枯叶组织在可见冠层投影中的像素比例。BiMSeg 通过排除背景限定统计范围，XGBoost 通过区分绿叶与枯叶确定其中计入分子的像素。冠层范围固定时，枯叶漏判与绿叶误判分别使估计值偏低和偏高；因此，枯叶 F1 为判断分类环节的可靠性提供了依据，而完整流程的图像级测量误差可通过模型估计值与人工参考枯叶率的逐图比较加以量化。